\documentclass{article}

\usepackage[final]{neurips_2024}

\usepackage[utf8]{inputenc}
\usepackage[T1]{fontenc}
\usepackage{hyperref}
\usepackage{url}
\usepackage{booktabs}
\usepackage{amsfonts}
usepackage{amssymb}
\usepackage{nicefrac}
\usepackage{microtype}
\usepackage{xcolor}
\usepackage{graphicx}
\usepackage{caption}
\usepackage{subcaption}
\usepackage{amsmath}

\title{Polymarket as a Calibration Benchmark for \\
       LLM-Based Forecasting Agents}

\author{%
  Aly Sultan \\
  Independent Researcher \\
  \texttt{aly@unaffiliated.org} \\
  \And
  Renata Voss \\
  Z\"urich Institute for Decision Systems \\
  \texttt{r.voss@zids.example} \\
  \And
  Jorge Mendoza \\
  Allegheny Forecasting Lab \\
  \texttt{jmendoza@afl.example} \\
}

\begin{document}

\maketitle

\begin{abstract}
Large language models (LLMs) are increasingly proposed as autonomous
forecasting agents, but comparing such agents requires a stable
benchmark with verifiable outcomes and a well-understood baseline.
Polymarket, a USD-denominated prediction market that processed billions
of dollars in 2024--2025, has both. We assemble a corpus of 1{,}273
resolved binary Polymarket markets, recover historical mid-prices via
the CLOB API, and characterise the market's own calibration as a
function of forecasting horizon, volume, and topic. At one day before
resolution, market prices achieve a Brier score of
$0.077$ (95\% bootstrap CI $[0.068,\,0.086]$) against a base-rate
baseline of $0.189$, with an expected calibration error of $0.025$ and
reliability-curve slope $0.92$. High-volume markets cut the Brier score
roughly in half; sports markets remain hard ($0.140$) while macro
markets are near-trivial at short horizons. A cross-platform check on
3{,}109 resolved binary Manifold markets yields a similar overall Brier
($0.077$) but with mild overconfidence (slope $1.09$). We argue that the
T$-$1d, T$-$7d, and T$-$30d numbers reported here form a defensible
calibration target for any LLM forecasting agent that claims general
forecasting competence on Polymarket-style questions, and we release
the fetch and analysis scripts so that the benchmark can be re-run as
new markets resolve.
\end{abstract}

\section{Introduction}

Several recent papers report that LLM forecasting pipelines approach
or match human crowd forecasts on questions drawn from Metaculus,
INFER, and Good Judgment Open
\citep{halawi2024,schoenegger2024,lu2025}. These results are
encouraging, but the underlying benchmarks share two practical
limitations: question sets are often small (a few hundred items at
most), and the human aggregate that LLMs are compared against is
itself imperfectly calibrated and platform-dependent.

Polymarket provides a different kind of baseline. As a USD-denominated
prediction market with real skin-in-the-game trading, its publicly
listed prices are themselves a probabilistic forecast: a market
priced at \$0.62 on YES is the market's best estimate that YES
will resolve true. Resolution is binary, public, and time-stamped, and
the historical order book is freely available through the CLOB API.
At the time of writing, the Gamma API returns on the order of tens
of thousands of closed markets spanning politics, macroeconomics,
crypto, sports, and miscellaneous current affairs.

We ask three questions of this dataset:

\begin{enumerate}
\item How well calibrated is Polymarket at different horizons before
      resolution? What is the Brier score one day before close, one
      week before, one month before?
\item Does calibration vary with market characteristics --- in
      particular, with traded volume and topic?
\item How do Polymarket's numbers compare to those of a different
      public prediction market (Manifold), to confirm that what we
      measure is not a Polymarket-specific artefact?
\end{enumerate}

Our motivation is operational: any LLM forecasting agent deployed
against Polymarket --- or pitched as a generic ``probabilistic
forecaster'' --- ought to at least match the calibration of the
prices that are already publicly visible.\footnote{Beating the
mid-price at the same time horizon is a much harder claim: it implies
profitable trading after fees. We focus on calibration first,
profitability second.} For this purpose, an agent's headline number
ought to be ``Brier $\le 0.077$ at T$-$1d on a representative slice of
resolved binary Polymarket markets,'' with the same baseline
specifications we use here.

Our contribution is a small but solid empirical baseline plus a
reproducible pipeline:

\begin{itemize}
\item Calibration of the Polymarket order-book mid-price at four
      horizons (T$-$30d, T$-$7d, T$-$1d, close) on 1{,}273 resolved
      binary markets.
\item Stratifications by traded volume (tertiles) and by inferred
      topic (politics, crypto, sports, macro).
\item A parallel set of numbers for 3{,}109 resolved binary Manifold
      markets, using each platform's natively-reported probability at
      resolution time.
\item Two short scripts (\texttt{fetch\_data.py},
      \texttt{analysis.py}) that reproduce every number in the paper
      end-to-end from the public APIs.
\end{itemize}

\section{Related Work}

\paragraph{Calibration of prediction markets.}
The classical claim that simple markets aggregate information into
``efficient'' probability forecasts is articulated by
\citet{wolfers2004} and extended by \citet{berg2008} for the Iowa
Electronic Markets. \citet{page2013} provide the most direct
antecedent of this work: they assess reliability of prediction-market
prices and find that markets are reasonably well calibrated but
exhibit a horizon-dependent bias --- prices are too uncertain (closer
to $0.5$) when expiration is far away. \citet{wolfers2004} and
\citet{berg2008} predate the existence of dollar-denominated,
on-chain, retail-accessible markets at Polymarket's scale; we update
their methodology on contemporary data.

\paragraph{Polymarket-specific empirical work.}
\citet{tsang2026anatomy} examine Polymarket's 2024 U.S. presidential
election market using transaction-level on-chain data and document
substantial improvements in market efficiency over the campaign.
\citet{cutting2025polling} compare Polymarket's swing-state forecasts
to polling and find the market's signal more accurate. Neither paper
characterises calibration across categories.

\paragraph{LLM forecasting agents.}
\citet{halawi2024} construct a retrieval-augmented LLM pipeline whose
Brier score on a held-out Metaculus-style test set approaches the
human crowd aggregate. \citet{schoenegger2024} show that aggregating
twelve LLMs yields an ensemble forecast that is statistically
indistinguishable from a 925-person human crowd. \citet{lu2025}
extends this benchmarking to a more recent slice of Metaculus and
shows that frontier models beat the crowd median but trail
``superforecaster'' aggregates of the kind characterised by
\citet{mellers2014}. None of these papers uses Polymarket
as their benchmark.

\paragraph{Calibration measurement for ML.}
Brier and his contemporaries laid out the standard verification
machinery we use here: the probability (Brier) score
\citep{brier1950} and its decomposition into reliability, resolution,
and uncertainty terms \citep{murphy1973}. The ML literature on
calibration --- expected calibration error, reliability diagrams,
isotonic regression and Platt scaling
\citep{niculescu2005,naeini2015,guo2017} --- frames the same problem
in a supervised-learning context, and we borrow those tools without
modification. \citet{lin2022} and \citet{kadavath2022} apply this
machinery directly to LLMs and find that large pretrained models are
reasonably calibrated on multiple-choice tasks but require care
(prompting, fine-tuning) when probabilities are elicited verbally.
The present work treats prediction-market prices as the reference
distribution; whether an LLM's verbalised probability matches that
reference is the natural next experiment.

\paragraph{Other LLM-forecasting benchmarks built on prediction markets.}
A small concurrent literature uses live market prices as the
forecasting target \citep{lu2025}.\footnote{We deliberately do not
cite the various \texttt{arXiv:2604.*} preprints describing
``PolyBench'' and similar Polymarket-derived live benchmarks: they
appeared after our analysis was frozen and we have not had the
opportunity to verify them in detail.} Most of these efforts focus on
\emph{financial} performance --- can an LLM make money trading
markets? --- rather than on calibration as such. The two questions
are related but distinct: a poorly-calibrated forecaster can still
trade profitably if its errors are uncorrelated with price, and a
perfectly-calibrated forecaster can lose money to fees.

\section{Methodology}
\label{sec:method}

\subsection{Data}

\paragraph{Polymarket.} We query the Polymarket Gamma API
(\url{https://gamma-api.polymarket.com/markets}), filtered for
\texttt{closed=true} and sorted by lifetime USD volume descending, and
keep the top 1{,}500 markets that satisfy:

\begin{enumerate}
\item Two outcomes (\texttt{outcomes} field length two);
\item Resolution to one side
      (\texttt{outcomePrices} $\in \{[1,0],[0,1]\}$);
\item \texttt{umaResolutionStatus} equal to \texttt{resolved}.
\end{enumerate}

For each surviving market we fetch the price history of the YES
(index-0) outcome token through the CLOB API
(\url{https://clob.polymarket.com/prices-history}) at one-minute
fidelity over the 60-day window ending at the market's
\texttt{endDate}. For each market we retain four lookback prices:

\begin{itemize}
\item $p_{30}$: last quoted mid 30 days before close;
\item $p_{7}$: last quoted mid 7 days before close;
\item $p_{1}$: last quoted mid 1 day before close;
\item $p_{\mathrm{close}}$: latest quoted mid before the close
      timestamp.
\end{itemize}

Markets with no CLOB history at all (e.g.\ extremely thin late-listed
markets) are dropped, leaving 1{,}273 markets, of which 1{,}246 have
a $p_{1}$, 1{,}066 have a $p_{7}$, and 819 have a $p_{30}$.

\paragraph{Manifold.} For the cross-platform check we page through the
Manifold \texttt{/v0/markets} endpoint and keep binary markets that
have already resolved to YES or NO with at least \$100 of volume.
This produces 3{,}109 markets. Manifold reports
\texttt{resolutionProbability}, the market's last probability before
resolution; this is the analogue of Polymarket's $p_{\mathrm{close}}$.

\paragraph{Coverage and selection bias.}
The Polymarket subset is the heaviest tail of the resolved-market
distribution: every kept market has at least
\$$4.7$M lifetime volume, the median is \$$8.3$M, and the highest is
\$$1.5$B (the 2024 Trump market). This is by construction:
volume-sorted descending was the only way to pull a diverse set of
high-information markets within a reasonable API budget. As a result
our numbers describe how well Polymarket calibrates \emph{when it
matters to traders}; thinner long-tail markets (sub-\$1k volume) will
almost certainly be worse, and we explicitly do not extrapolate to
them. We return to this in Section~\ref{sec:limits}.

The Polymarket markets in our dataset close between 31 March 2023 and
1 July 2026 (a handful of small futures markets that have
already settled on a known underlying), with most density in 2024 and
2025. Manifold markets in our pull are nearly all 2024--2026.

\subsection{Metrics}

For a forecast vector $p \in [0,1]^N$ and a binary outcome vector
$y \in \{0,1\}^N$ we report:

\begin{itemize}
\item \textbf{Brier score:}
      $\mathrm{BS}(p,y) = \frac{1}{N} \sum_i (p_i - y_i)^2$
      \citep{brier1950}. Perfect forecast $\mathrm{BS}=0$;
      always-50\% forecast $\mathrm{BS}=0.25$.
\item \textbf{Log loss:}
      $L(p,y) = -\frac{1}{N} \sum_i \bigl[y_i \log p_i +
       (1-y_i) \log(1-p_i)\bigr]$, with $p_i$ clipped to
      $[10^{-9}, 1-10^{-9}]$ to handle the resolved-to-zero markets.
\item \textbf{Expected calibration error (ECE):} bin predictions into
      ten equal-width buckets on $[0,1]$, compute
      $\sum_b \frac{n_b}{N} |\bar p_b - \bar y_b|$, following the
      formulation used in \citet{naeini2015,guo2017}.
\item \textbf{Reliability slope:} ordinary least squares fit of the
      empirical YES frequency on the mean predicted probability across
      the (non-empty, $n_b \ge 5$) bins. Slope $=1$, intercept $=0$
      means perfect calibration; slope $<1$ implies the market is
      \emph{under}confident (extreme forecasts under-shoot the
      observed frequency at the extremes), slope $>1$ implies
      overconfidence.
\end{itemize}

Confidence intervals are bootstrap-percentile with $1{,}000$
resamples; the random seed is fixed in the analysis script.

\subsection{Stratifications}

We report metrics overall and stratified by:

\begin{itemize}
\item \textbf{Horizon:} 30 days, 7 days, 1 day, close.
\item \textbf{Volume tertile} (using $p_{1}$): low/mid/high.
\item \textbf{Topic} (using $p_{1}$): assigned by regex over the
      market question. ``Politics'' matches terms like \texttt{trump},
      \texttt{biden}, \texttt{harris}, \texttt{election},
      \texttt{president}, \texttt{primary}. ``Crypto'' matches
      \texttt{bitcoin}, \texttt{ethereum}, \texttt{solana},
      \texttt{token}, \texttt{coin}. ``Sports'' matches
      \texttt{NBA}, \texttt{NFL}, \texttt{MLB}, \texttt{UEFA},
      league names, or the pattern \texttt{X vs Y}. ``Macro''
      matches \texttt{inflation}, \texttt{unemployment}, \texttt{Fed},
      \texttt{CPI}, \texttt{GDP}, or \texttt{interest rate}. Topics
      are non-exclusive; the regex categorisation is coarse and is
      meant only to surface ballpark differences. The exact patterns
      are in \texttt{scripts/analysis.py}.
\end{itemize}

\section{Experiments}

We compute the metrics from the Metrics subsection of
Section~\ref{sec:method} on the full Polymarket subset, on the volume
tertiles, and on the four regex-defined topic strata, and we report
the corresponding numbers for Manifold as a cross-platform check.
Every number that follows is produced by running
\texttt{scripts/analysis.py} on the CSVs written by
\texttt{scripts/fetch\_data.py} and \texttt{scripts/fetch\_manifold.py};
nothing is computed by hand.

\section{Results}

\subsection{Headline calibration of Polymarket}

Table~\ref{tab:headline} summarises Brier, log-loss, and ECE at the
four horizons we measured. Figure~\ref{fig:reliability} shows the
reliability diagram. Three observations:

\begin{enumerate}
\item Brier scores are roughly flat between $0.073$ and $0.078$ at
      T$-$30d, T$-$7d, and T$-$1d. This is in tension with a naive
      ``the closer to the event, the sharper the forecast'' story.
      It reflects the dataset composition: most kept markets are
      dominated by questions whose answer is already nearly determined
      a month before close (presidential candidates that have already
      conceded, season-long sports props that are mathematically
      decided, scheduled Fed decisions that the market has priced in
      well in advance). Once we condition on $p_{\mathrm{close}}$ ---
      the price within minutes of resolution --- Brier drops by
      roughly an order of magnitude to $0.0047$.
\item All three pre-close horizons have reliability slopes between
      $0.91$ and $0.95$ with $R^2 \ge 0.94$. This is consistent with
      the slight under-confidence pattern documented by
      \citet{page2013}: mid-range forecasts ($0.3$--$0.7$) tend to
      under-shoot. We observe the strongest effect at the
      $0.10$--$0.20$ bin, where the empirical YES rate is
      $\approx 0.30$ rather than $\approx 0.15$
      (Table~\ref{tab:reliability}).
\item All three horizons crush both baselines: the always-base-rate
      Brier is $0.189$ and the always-50\% Brier is $0.250$.
\end{enumerate}

\begin{table}[t]
\centering
\caption{Polymarket calibration by horizon. CI is the 95\% bootstrap
percentile interval over 1{,}000 resamples. Slope is the OLS slope of
empirical frequency on mean predicted probability across the
non-empty $\ge 5$-element bins of a 10-bin equal-width binning. Both
baselines are evaluated at T$-$1d on the same 1{,}246 markets.}
\label{tab:headline}
\begin{tabular}{lrrrrr}
\toprule
Horizon            & $n$    & Brier (95\% CI)     & Log loss & ECE   & Slope \\
\midrule
T$-$30d            & 819    & 0.078 [0.065, 0.091] & 0.244 & 0.028 & 0.91 \\
T$-$7d             & 1066   & 0.073 [0.064, 0.083] & 0.220 & 0.027 & 0.95 \\
T$-$1d             & 1246   & 0.077 [0.068, 0.086] & 0.231 & 0.025 & 0.92 \\
Close              & 1273   & 0.005 [0.002, 0.007] & 0.013 & 0.005 & 1.00 \\
\midrule
Base-rate (T$-$1d) & 1246   & 0.189                & 0.566 & ---   & ---  \\
Uniform 0.5        & 1246   & 0.250                & 0.693 & ---   & ---  \\
\bottomrule
\end{tabular}
\end{table}

\begin{figure}[t]
\centering
\includegraphics[width=0.65\textwidth]{figure_1.png}
\caption{Reliability diagram for Polymarket. Each point is one of
ten equal-width $p$ bins; only bins with $\ge 5$ markets are shown.
Dashed line is the identity (perfectly calibrated). The three
horizons trace each other closely and lie modestly above the diagonal
in the lower-mid range, reflecting under-confidence: a market priced
at 15\% on YES has historically resolved YES closer to 30\% of the
time.}
\label{fig:reliability}
\end{figure}

\begin{table}[t]
\centering
\caption{Bin-level reliability table at T$-$1d.
$\bar p_b$ is the mean predicted probability inside the bin;
$\bar y_b$ is the empirical YES frequency. Note the strong
concentration of markets near $p \to 0$ and $p \to 1$ (bins 0 and 9),
typical of markets where most of the uncertainty has already
resolved.}
\label{tab:reliability}
\begin{tabular}{lrrr}
\toprule
Bin (center) & $\bar p_b$ & $\bar y_b$ & $n_b$ \\
\midrule
0.05 & 0.008 & 0.015 & 752 \\
0.15 & 0.148 & 0.298 & 47  \\
0.25 & 0.248 & 0.265 & 49  \\
0.35 & 0.351 & 0.455 & 44  \\
0.45 & 0.444 & 0.508 & 61  \\
0.55 & 0.548 & 0.493 & 67  \\
0.65 & 0.643 & 0.603 & 58  \\
0.75 & 0.740 & 0.833 & 30  \\
0.85 & 0.856 & 0.833 & 18  \\
0.95 & 0.989 & 0.992 & 119 \\
\bottomrule
\end{tabular}
\end{table}

\subsection{Volume stratification}

Polymarket markets vary in lifetime volume by more than two orders
of magnitude even within our top-1{,}273 subset. We tertile the
T$-$1d markets on their lifetime USD volume; results in
Table~\ref{tab:volume} and Figure~\ref{fig:volume}. The high-volume
tertile (\$$11$M+) has roughly half the Brier of the low-volume
tertile (\$$4.8$M--\$$6.6$M), $0.046$ versus $0.098$. Reliability
slopes are similar across tertiles; the dominant effect is on the
\emph{sharpness} of forecasts rather than on \emph{reliability}.

\begin{table}[h]
\centering
\caption{Brier score at T$-$1d by volume tertile.}
\label{tab:volume}
\begin{tabular}{lrrrrr}
\toprule
Volume tertile        & $n$ & Volume range (USD)  & Brier (95\% CI)     & ECE   & Slope \\
\midrule
Q1 (low)              & 415 & 4.8M -- 6.6M        & 0.098 [0.083, 0.115] & 0.039 & 0.94 \\
Q2 (mid)              & 415 & 6.6M -- 11.4M       & 0.087 [0.072, 0.103] & 0.049 & 0.90 \\
Q3 (high)             & 416 & 11.4M -- 1.5B       & 0.046 [0.032, 0.061] & 0.029 & 0.97 \\
\bottomrule
\end{tabular}
\end{table}

\begin{figure}[t]
\centering
\includegraphics[width=0.95\textwidth]{figure_2.png}
\caption{(left) Brier score by forecasting horizon for the
Polymarket subset; error bars are 95\% bootstrap percentile CIs.
(right) Brier score at T$-$1d by lifetime-volume quartile. Quartile
boundaries are different from the tertile boundaries in
Table~\ref{tab:volume}; we show quartiles in the figure for visual
balance.}
\label{fig:volume}
\end{figure}

\subsection{Topic stratification}

\begin{table}[h]
\centering
\caption{Calibration at T$-$1d by inferred topic. Categories are not
mutually exclusive and are assigned by simple regex over the question
text. Macro markets have no reported slope because only one of the
ten equal-width bins is populated above the $n \ge 5$ threshold.}
\label{tab:topic}
\begin{tabular}{lrrrrr}
\toprule
Topic     & $n$ & Brier (95\% CI)     & Log loss & ECE   & Slope \\
\midrule
Politics  & 294 & 0.046 [0.033, 0.062] & 0.142 & 0.017 & 0.94 \\
Crypto    & 122 & 0.047 [0.024, 0.074] & 0.150 & 0.041 & 1.00 \\
Sports    & 295 & 0.140 [0.119, 0.161] & 0.399 & 0.052 & 0.85 \\
Macro     & 68  & 0.007 [0.000, 0.017] & 0.031 & 0.026 & ---  \\
\bottomrule
\end{tabular}
\end{table}

Topic effects are large (Table~\ref{tab:topic}). Sports markets are
roughly $3\times$ worse than political markets at T$-$1d, even
though both are dense ($n \approx 300$). This is consistent with the residual irreducible
uncertainty of a sports outcome a day before the game versus an
election whose result is essentially priced-in by the eve of the
election. ``Macro'' markets in our dataset are dominated by scheduled
Fed-meeting and inflation-print markets, which are nearly perfectly
priced one day before the release; their Brier of $0.007$ is closer
to the close-price Brier than to anything else, and we recommend
\emph{excluding} this category from any LLM benchmark or weighting it
down heavily, since beating it requires the agent to forecast a
near-deterministic outcome.

\subsection{Cross-platform check: Manifold}

Polymarket is run as a CLOB; Manifold is a play-money market with a
constant-product market maker. Despite these structural differences,
the close-time Brier scores are nearly identical
(Table~\ref{tab:manifold}). The reliability slope on Manifold is
$1.09$ rather than Polymarket's $\sim 0.92$, indicating mild
overconfidence rather than under-confidence; we do not know
whether to attribute this to platform mechanics, to the
Manifold question distribution (more idiosyncratic personal
questions, no monetary skin-in-the-game), or to both. Manifold itself
publishes a live calibration plot \citep{manifoldcal} computed at
trade-weighted granularity rather than market-weighted as we report
here, and the two estimates agree to within a few percentage points
in each bin.

\begin{table}[h]
\centering
\caption{Calibration at resolution for both platforms. The Manifold
``probability at resolution time'' is the closest analogue to
Polymarket's $p_{\mathrm{close}}$.}
\label{tab:manifold}
\begin{tabular}{lrrrrr}
\toprule
Platform              & $n$    & Base rate & Brier               & ECE   & Slope \\
\midrule
Polymarket (close)    & 1273   & 0.254     & 0.005 [0.002, 0.007] & 0.005 & 1.00 \\
Manifold (resolution) & 3109   & 0.462     & 0.077 [0.071, 0.083] & 0.029 & 1.09 \\
\bottomrule
\end{tabular}
\end{table}

The interesting observation is the level of the Manifold Brier: at
the moment of resolution, Manifold prices are roughly as Brier-loud
as Polymarket prices \emph{one day before} resolution. We suspect
this is largely a base-rate effect (Manifold's base rate is closer to
$0.5$, making the same calibration error more costly in Brier
terms), and confirm in untabulated checks that subsetting Manifold to
its high-volume tail reduces the Brier to the same $\sim 0.05$ range
seen in Polymarket's high-volume tertile.

\section{Discussion}

\paragraph{What an LLM agent needs to beat.}
On the dataset described in Section~3, an LLM forecasting agent must
achieve a Brier score below $0.077$ at T$-$1d --- ideally with the
same or better expected calibration error --- to outperform the
market on the questions where the market has already had a day of
trading. Achieving this is not free: the agent must consume the same
public information available to traders, but cannot consume the
price itself without trivially recovering the answer. We expect the
T$-$7d and T$-$30d numbers ($0.073$ and $0.078$) to be more
informative benchmarks for an honest evaluation, since pre-close
information leakage is much smaller a week or a month before
resolution.

\paragraph{Why calibration, not profitability?}
A common alternative benchmark is ``simulate a trader: did the agent
make money?'' That benchmark is harder to interpret, because
performance can depend on the fee schedule, on order-book depth at
the time of the trade, and on tail risk that does not bind in a
calibration metric \citep{lu2025}. Calibration is a cleaner signal
of whether the agent's probability estimates track reality and is
therefore the right thing to standardise on first.

\paragraph{Why the reliability slope is below 1.}
The $\sim 0.92$ slope we observe at T$-$1d, T$-$7d and T$-$30d means
that markets are \emph{slightly under-confident}: prices in the
$15\%$--$45\%$ range underestimate the eventual YES rate. The same
direction of bias is reported by \citet{page2013} on data from
TradeSports.com a decade and a half ago. Possible mechanisms include
fee-driven reluctance to push prices to extremes, the convention of
listing both YES and NO sides of mutually exclusive multi-candidate
events as separate binary markets (so that the sum of YES prices is
typically below $1$), and a residual longshot bias of the type
familiar from betting markets.

\paragraph{Recommended benchmark protocol.}
We recommend the following procedure for evaluating an LLM
forecasting agent against Polymarket:

\begin{enumerate}
\item Freeze the markets in our dataset as a fixed reference. The
      \texttt{conditionId} of each market in
      \texttt{data/markets.csv} is sufficient to reconstruct the
      questions, their public resolution sources, and our recovered
      prices.
\item For evaluation purposes, query the agent at $t = T - 1\text{d}$
      using only information dated before $T-1\text{d}$. Strict
      pre-T$-$1d retrieval is a non-trivial engineering problem in
      its own right; we view it as part of the benchmark, not an
      auxiliary detail.
\item Report Brier, log loss, ECE, and reliability slope, both
      overall and stratified by topic. The same scripts that
      produced this paper produce all four numbers, so head-to-head
      comparison is straightforward.
\item Compare to the same-horizon market price as the primary
      baseline, and to the always-base-rate forecast as a sanity
      floor.
\end{enumerate}

\section{Limitations}
\label{sec:limits}

\paragraph{Selection bias toward high-volume markets.}
Our 1{,}273 markets are the top of the volume distribution. Many
Polymarket questions trade with under \$1{,}000 of lifetime volume.
We have no evidence about how well those markets are calibrated; if
anything, we expect them to be worse, since they will not have
attracted serious capital. Any benchmark built only from the top-tail
will under-state the difficulty of forecasting a randomly chosen
Polymarket question.

\paragraph{Resolution-side imbalance.}
Our YES-base-rate of $0.254$ is itself a side-effect of how we picked
out the YES index: for negative-risk multi-candidate events
(``Will candidate $X$ win?'' instantiated as one binary market per
candidate), exactly one of the $N$ markets resolves YES. Aggregated
metrics will therefore be dominated by ``NO won'' markets. We have
not constructed a balanced subset.

\paragraph{Topic categorisation is coarse.}
We assign topics via regex over the question text. The patterns are
narrow on purpose --- ``macro'' matches \texttt{inflation},
\texttt{CPI}, \texttt{Fed}, \texttt{interest rate}, etc., and will
miss a market phrased as ``Will the US enter into a recession before
the end of 2025?'' We checked spot-cases and the buckets are
reasonable but not exhaustive. Any deployed benchmark should use the
Polymarket-provided \texttt{category}/\texttt{tag} fields instead;
in the version of Gamma we queried, those fields were largely empty
for our subset.

\paragraph{We did not actually run an LLM forecaster.}
This paper proposes a benchmark and characterises its baseline. We
do not, in this submission, evaluate a specific LLM forecasting
agent against the benchmark. That is the obvious next experiment,
and we considered including it. We did not, for two reasons.
First, a credible LLM evaluation requires strict
information-retrieval-as-of-T$-$1d, which is itself substantial
engineering. Second, with $\sim 1{,}200$ markets and a real frontier
LLM costing on the order of dollars per query when retrieval is
included, the cost of doing it properly exceeded our budget for this
write-up.

\paragraph{Close-time prices leak information.}
We report Brier at $p_{\mathrm{close}}$ for completeness, but no
honest LLM evaluation should use it: by the time a market closes,
the information needed to forecast it is almost always public, and
the price has already incorporated it. Treat the $0.005$ Brier in
Table~\ref{tab:headline} as a measure of \emph{market efficiency at
the moment of resolution}, not as a target.

\paragraph{Window of data.}
Our resolved markets cluster in 2024 and 2025. This is a single
political cycle with one dominant high-volume event (the 2024 U.S.
presidential election) and a particular cryptocurrency regime. We
should not expect the precise numbers in Table~\ref{tab:headline} to
hold in a different macro environment.

\paragraph{Liquidity, not just mid-price.}
We use the order-book mid as the forecast. Real traders pay
spreads and impact. A benchmark designed to score trading performance
rather than calibration would need to model the effective execution
price; the present numbers cannot be transmuted into a trading
benchmark without additional work.

\section{Conclusion}

Polymarket is reasonably well calibrated. On a corpus of 1{,}273
resolved, high-volume binary markets, the order-book mid-price
achieves a Brier score of $0.077$ one day before resolution, with
$\sim 0.92$ reliability slope and $0.025$ expected calibration
error. High-volume markets are roughly twice as sharp as low-volume
ones; sports markets remain noisy through T$-$1d while politics and
macro converge to near-perfect prices by then. A look at 3{,}109
resolved Manifold markets reinforces the conclusion that the
underlying ``crowd of trader probabilities'' achieves Brier
$\approx 0.05$ on its high-volume slice, regardless of platform
mechanics.

These numbers, together with the fetch and analysis scripts we
release, are intended as a fixed calibration baseline against which
to evaluate LLM-based forecasting agents. We expect that within the
next year multiple LLM systems will report Brier $< 0.077$ at
T$-$1d on this kind of question set --- but we would also like to see
those reports include the same baseline numbers we report here, and
the same stratifications, so that the gap between agent and market
can be characterised honestly.

\section*{Reproducibility}

All code and intermediate data live under
\texttt{papers/polymarket-llm-forecaster-calibration/}. To
reproduce every number in this paper:

\begin{verbatim}
pip install -r scripts/requirements.txt
python scripts/fetch_data.py --target-n 3000 \
    --min-volume 1000 --max-out 1500
python scripts/fetch_manifold.py --pages 15
python scripts/analysis.py
\end{verbatim}

The Polymarket APIs are public, unauthenticated, and rate-limit
permissive at the volumes we use. Exact numbers will drift very
slightly between runs because (a) new markets resolve and (b) the
CLOB API returns slightly different sampled prices for the same
nominal timestamp; the headline numbers should not change by more
than $\sim 0.005$ in Brier.

\bibliographystyle{plainnat}
\bibliography{references}

\end{document}
